\chapter*{Conclusion générale et perspectives}
\addcontentsline{toc}{chapter}{Conclusion générale et perspectives}
\markboth{Conclusion générale et perspectives}{}

Ce stage avait pour objectif de permettre à un opérateur de commander le robot AMR-X
depuis le dashboard web du projet et d'y suivre fidèlement sa position et le résultat de
sa navigation. Au départ, le dashboard, la pile Nav2 et la simulation Gazebo
fonctionnaient séparément.

Pour les relier, une passerelle ROS~2 dédiée a été conçue et réalisée. Elle reçoit les
destinations par un service acquitté, vérifie chaque demande, dialogue avec Nav2 et publie
un état de navigation unique et versionné. Côté dashboard, la page \emph{Live map} permet
désormais de choisir une destination sur la carte ou par saisie, de l'envoyer, de suivre
la navigation et de l'annuler, tout en bloquant les commandes lorsque les données ne sont
plus fiables. Un script de lancement rend la chaîne complète reproductible en une seule
commande.

La validation en simulation a montré que la destination est acquittée immédiatement, que
l'état affiché suit l'exécution réelle de Nav2, que les demandes invalides ou concurrentes
sont refusées et que l'annulation arrête le robot. Elle a surtout mis en évidence un écart
de 4,1~m entre la position affichée et la position réelle du robot, dû à la différence
d'origine entre la carte SLAM et le monde Gazebo. Après mesure et configuration de ce
décalage, la position affichée dans le dashboard correspond à celle du robot dans Gazebo
à moins de 6~cm près. Ce résultat illustre l'importance de comparer l'interface à une
vérité terrain plutôt que de se fier à la seule cohérence interne du système.

Plusieurs perspectives se dégagent de ce travail :
\begin{itemize}
  \item terminer les scénarios restant à valider depuis l'interface (obstacle,
    destination inaccessible, coupure de rosbridge, mode démonstration, changement
    d'adresse) ;
  \item déterminer automatiquement le décalage entre la carte et le monde, par exemple en
    l'enregistrant avec la carte lors du SLAM, au lieu de le configurer à la main ;
  \item enregistrer les navigations lancées depuis la \emph{Live map} dans la base de
    données, afin de les relier à la gestion des missions ;
  \item étendre la passerelle à plusieurs robots et à l'arbitrage entre sources de
    commande ;
  \item sécuriser l'accès à rosbridge avant tout usage en dehors d'un poste de simulation,
    puis préparer le passage sur le robot physique.
\end{itemize}
